% chapters/ch11_debugging_lab.tex — Chapter 11: Hardware Debugging & Lab Skills
% Scope (PLAN.md): oscilloscope (triggering, probing UART/SPI/CAN, reading the
% waveform), logic analyzer vs scope, JTAG/SWD internals, ST-Link vs J-Link,
% breakpoints vs watchpoints, board bring-up sequence, reading schematics,
% ESD & lab safety. Budget 30 pp, mix 22 Q&A / 4 code / 4 concept.
% All snippets verified gcc -Wall -Wextra -std=c11 (zero warnings).

\chapter{Hardware Debugging \& Lab Skills}
\label{ch:debugging-lab}

\begin{partnote}
\textbf{Scope} --- the oscilloscope (triggering, probing, reading UART/SPI/CAN
on screen) $\rightarrow$ logic analyser vs scope $\rightarrow$ JTAG/SWD
internals $\rightarrow$ ST-Link vs J-Link $\rightarrow$ breakpoints vs
watchpoints $\rightarrow$ the board bring-up sequence $\rightarrow$ reading
schematics $\rightarrow$ ESD \& lab safety. The whole book keeps invoking
``scope it at the receiver''; this chapter makes the instrument-led method
explicit. It is the candidate's daily reality --- the lab instruments and
ST-Link/J-Link on the resume.
\end{partnote}

The difference between an engineer who \emph{guesses} and one who \emph{knows} is
the lab bench. When firmware and hardware disagree, the scope, logic analyser,
and debugger turn ``it doesn't work'' into a measured fact. These questions probe
whether you've actually held the probes.

%=====================================================================
\section{Primer}
%=====================================================================

\subsection{Oscilloscope vs logic analyser}

An \textbf{oscilloscope} shows \emph{analog} voltage versus time on a few
channels with fine vertical/timing resolution --- it tells you signal
\emph{integrity}: edge shape, ringing, overshoot, noise, levels, rise time. A
\textbf{logic analyser} shows \emph{many digital} channels as 1s and 0s with
\emph{protocol decode} --- it tells you signal \emph{content}: the bytes on a
UART/SPI/I2C/CAN bus, timing relationships across many lines. Rule of thumb: use
the scope when you suspect the \emph{electrical} (is the edge clean? is the level
right? is it terminated?), the logic analyser when you suspect the \emph{logical}
(are the right bytes being sent? is CS framing correct?). Many bring-up problems
need both.

\subsection{Triggering --- the key scope skill}

A scope is only useful if it captures the \emph{right} moment. The
\textbf{trigger} defines that: \textbf{edge} trigger (fire on a rising/falling
edge at a voltage threshold) is the default; \textbf{pulse-width} trigger
catches runts or glitches of a given width; \textbf{serial/protocol} trigger
fires on a specific byte/address/condition on UART/SPI/I2C/CAN; \textbf{single
shot} captures a one-time event. Without a stable trigger the waveform ``runs''
and you see nothing useful. The classic technique: toggle a spare GPIO in your
code at the event of interest and trigger the scope on that GPIO --- turning a
software moment into a hardware trigger (used for measuring interrupt latency in
Chapter 4).

\subsection{JTAG, SWD, and the debug probes}

\textbf{JTAG} is the classic 4--5-wire (TCK/TMS/TDI/TDO/\,nTRST) boundary-scan
and debug interface; \textbf{SWD} (Serial Wire Debug) is ARM's 2-wire
(SWCLK/SWDIO) alternative that does the same core debug over fewer pins ---
preferred on pin-constrained MCUs. Both let a host \textbf{debug probe} halt the
core, read/write registers and memory, set breakpoints, and flash the chip. An
\textbf{ST-Link} is ST's probe (bundled on every Nucleo/Discovery board, great
value for STM32); a \textbf{J-Link} (Segger) is a higher-end probe with broad
device support, faster flashing, and features like unlimited flash breakpoints
and RTT. The candidate uses both (resume), choosing by target and feature need.

\subsection{The board bring-up sequence}

A new board comes up in a fixed order, each step a precondition for the next:
\textbf{power} (verify every rail's voltage and ripple with a meter/scope ---
nothing else matters if a rail is wrong) $\rightarrow$ \textbf{clocks} (confirm
the crystal/oscillator is running at the right frequency) $\rightarrow$
\textbf{reset} (the reset line releases cleanly, no glitches) $\rightarrow$
\textbf{debug access} (the probe can connect over JTAG/SWD and halt the core)
$\rightarrow$ \textbf{peripherals} (bring up GPIO/UART/buses one at a time). Skip
a step and you chase ghosts --- a firmware bug that's really a sagging rail or a
dead crystal. ``Power $\rightarrow$ clocks $\rightarrow$ reset $\rightarrow$ debug
$\rightarrow$ peripherals'' is the litany.

%=====================================================================
\section{Q\&A Bank}
%=====================================================================

\tierbanner{Tier 1 --- Screening (phone / HR-technical)}%
{Two to four sentences.}

\question{Q1.1}{Oscilloscope vs logic analyser --- when each?}
A scope shows analog voltage vs time on a few channels --- use it for signal
integrity (edges, ringing, levels, noise). A logic analyser shows many digital
channels with protocol decode --- use it to see the actual bytes/timing on a bus.
Electrical problem $\rightarrow$ scope; logical/content problem $\rightarrow$
logic analyser.

\question{Q1.2}{What is triggering and why does it matter?}
The trigger tells the scope \emph{when} to capture, so a repeating or one-shot
event is shown stably instead of an unstable blur. Without a proper trigger
(edge, pulse-width, protocol, single-shot) you can't see the event you care
about.

\question{Q1.3}{JTAG vs SWD?}
Both are ARM debug interfaces; JTAG uses 4--5 wires (TCK/TMS/TDI/TDO/nTRST) and
supports boundary scan, while SWD uses just 2 (SWCLK/SWDIO) for the same core
debug --- preferred when pins are scarce. Both let a probe halt the core, access
memory/registers, and flash the device.

\question{Q1.4}{ST-Link vs J-Link?}
ST-Link is ST's debug probe, bundled on STM32 dev boards --- cheap and ideal for
STM32. J-Link (Segger) is a higher-end probe with wide device support, faster
programming, and features like unlimited breakpoints and RTT logging. You pick by
target and the features you need.

\question{Q1.5}{Breakpoint vs watchpoint?}
A breakpoint halts the core when execution reaches a specific \emph{address}
(instruction); a watchpoint (data breakpoint) halts when a specific \emph{memory
location} is read or written, regardless of where in the code. Watchpoints are
how you catch \emph{who} corrupted a variable.

\question{Q1.6}{What's the first thing you check bringing up a new board?}
Power --- measure every supply rail for correct voltage and acceptable ripple
before anything else. A wrong or noisy rail makes every higher-level symptom
meaningless, so power is always step one.

\question{Q1.7}{Why probe a signal at the receiver, not the transmitter?}
Because the receiver sees what actually arrived after the cable, connectors, and
termination have had their effect --- reflections, attenuation, and noise show up
there. A signal can look clean at the driver and be corrupted at the far end,
which is where the failure actually occurs.

\question{Q1.8}{What is ESD and why care in the lab?}
Electrostatic discharge --- a static voltage spike (from your body, clothing)
that can instantly damage or, worse, \emph{latently} weaken a semiconductor. You
mitigate it with a grounded wrist strap, ESD mat, and careful handling, because a
zapped part may fail now or intermittently weeks later.

\question{Q1.9}{What does a DMM tell you that a scope doesn't, and vice versa?}
A DMM gives precise steady-state values (DC voltage, resistance, continuity,
current) but not fast dynamics. A scope shows the \emph{time-varying} behaviour
(edges, glitches, ripple, frequency) but with less DC precision. You use the DMM
for ``is this rail 3.3\,V and is this trace continuous?'' and the scope for ``what
does the edge look like?''

\question{Q1.10}{Why use $10\times$ probe attenuation?}
A $10\times$ probe presents higher input impedance and lower capacitance to the
circuit, so it loads (disturbs) the signal less and has more bandwidth --- at the
cost of $10\times$ smaller displayed amplitude (the scope compensates). For
fast/high-impedance nodes a $1\times$ probe's loading can distort or even stop the
signal.

\tierbanner{Tier 2 --- Core technical (panel round)}%
{A short paragraph.}

\question{Q2.1}{How do you read a UART byte on a scope/logic analyser?}
Trigger on the falling edge of the start bit (line idles high). Then sample at
the bit centres spaced by the baud period: the start bit (0), eight data bits
\emph{LSB first}, optional parity, and the stop bit (1). Decode the data bits
into the byte (a logic analyser does this automatically with protocol decode).
Confirm the bit width matches the configured baud --- a width mismatch is a
baud/clock error (Chapter 5). This is exactly the decode in the Coding Gym.

\question{Q2.2}{How do you measure interrupt latency on the bench?}
Toggle a GPIO high the instant the triggering event occurs (or in hardware) and
toggle another (or the same) GPIO at the very start of the ISR; trigger the scope
on the first edge and measure the time to the second. That interval is the
hardware-plus-software latency. The same GPIO-toggle trick measures task timing,
loop period, and jitter --- turning software events into scope-visible edges
(Chapter 4). For pure software timing you can also use the \code{DWT} cycle
counter.

\question{Q2.3}{Walk the board bring-up sequence and what tool you use at each
step.}
\textbf{Power}: DMM/scope on every rail (voltage + ripple). \textbf{Clocks}:
scope on the crystal/oscillator (or a clock-out pin) to confirm frequency.
\textbf{Reset}: scope the reset line for a clean release, no glitches/brownout.
\textbf{Debug}: connect the ST-Link/J-Link over SWD/JTAG and confirm the core
halts and the IDCODE reads. \textbf{Peripherals}: enable one at a time (GPIO
blink first, then UART console, then the buses), scoping/analysing each. Each step
gates the next, so a failure is localised rather than mysterious.

\question{Q2.4}{How does SWD physically work over two wires?}
SWDIO is bidirectional data and SWCLK is the clock; the host drives a packet ---
a request (read/write, AP/DP, address, parity), then a turnaround, an ACK from
the target, then the 32-bit data phase with its own parity. The single data line
reverses direction (turnaround cycles) between host and target. It accesses the
core's debug registers (the Debug Access Port) to halt, single-step, read/write
memory, and program flash --- all the JTAG debug functions over two pins. The
parity protects each phase (Coding Gym).

\question{Q2.5}{When do you reach for a logic analyser over a scope on a bus?}
When the electrical levels look fine but the \emph{data} is wrong, or when you
need to see many lines together (SPI's SCLK/MOSI/MISO/CS, or a parallel bus) and
decode the protocol. The logic analyser captures long sequences and shows
decoded bytes, addresses, ACK/NACK, etc., which a 2--4 channel scope can't
practically do. You'd still grab the scope if you suspect the levels/edges
themselves (ringing, slow rise, marginal thresholds).

\question{Q2.6}{What is a watchpoint good for, with an example?}
A watchpoint halts when a chosen variable/address is accessed, so it catches
\emph{memory corruption} without knowing where the offending code is: set a
write-watchpoint on the variable that mysteriously changes, run, and the debugger
stops exactly at the instruction that wrote it --- revealing a stray pointer,
buffer overrun, or stack overflow (Chapter 2). It's the fastest way to answer
``who is clobbering this?'' that a breakpoint (which needs you to already suspect
the code location) can't.

\question{Q2.7}{How do you read a schematic to debug a signal?}
Trace the net of interest: find the source (which pin drives it), follow it
through series resistors, pull-ups/downs, level shifters, and connectors to the
destination, noting reference designators so you can probe the physical points.
Check the power net feeding each part and the decoupling. The schematic tells you
\emph{what the signal should be and where to measure it}; the scope tells you what
it actually is. Net names and designators are the bridge between the drawing and
the board.

\question{Q2.8}{What is RTT and why is it useful?}
Segger's Real-Time Transfer is a way to stream \code{printf}-style logging out
through the debug probe (J-Link) via a RAM buffer the host reads over SWD ---
with almost no target overhead and \emph{no UART pin}. It's far faster and less
intrusive than a UART console and doesn't perturb timing the way a blocking
\code{printf} does, so you can log from time-critical code. It's a key reason to
pick a J-Link for serious debugging.

\tierbanner{Tier 3 --- Trap \& depth (principal-engineer level)}%
{A paragraph plus the trap.}

\question{Q3.1}{A signal looks perfect on the scope at the driver but the
receiver mis-reads it. Where's the fault, and how do you prove it?}
The corruption happens \emph{in transit} --- between where you probed and where
it's sampled --- so probing at the driver is the mistake. Likely causes:
reflections from missing/incorrect termination (rings worst at the far end),
attenuation/slow edges over a long or capacitive trace, crosstalk from an
adjacent aggressor, or a ground/common-mode shift between the two ends. Proof:
move the probe to the \emph{receiver} input pin and capture there; you'll see the
ringing/overshoot/slow edge that the driver-side view hid. \trap the trap is
trusting the clean driver-side capture and blaming firmware; the disciplined move
is always to scope at the point of \emph{sampling}, which is the recurring method
across the UART/SPI/CAN chapters and the resume's oscilloscope root-cause work.

\question{Q3.2}{The debugger ``works'' but the firmware behaves differently when
you run it standalone from reset. Enumerate causes and how the debugger hides
them.}
The debugger is a privileged, altered environment: it may set the stack pointer,
clock, and \code{VTOR} for you (so a wrong vector base or uninitialised clock
only bites from cold reset --- Chapters 2--3); it halts the core, which stops the
watchdog from firing (hiding a watchdog-reset loop) and serialises execution
(hiding race/timing bugs that need full speed); it may keep clocks/power-domains
active that low-power code would gate; and loading via the debugger initialises
RAM that a real reset leaves as garbage (hiding \code{.bss}/\code{.data} startup
bugs). \trap the trap is treating ``runs under the debugger'' as ``works''; it's a
different environment. The method is to reproduce from a true power-on reset and
instrument early --- toggle a GPIO at each startup milestone, check rails/clock/
\code{VTOR}/watchdog status --- rather than only stepping in the IDE.

\question{Q3.3}{How do you debug an intermittent fault that you can't reproduce
on demand?}
Make it observable and captured: set up a \textbf{trigger} on the failure
signature (a protocol error byte, a glitch of a certain width, a GPIO you toggle
in the fault path) and let the scope/logic analyser \emph{single-shot} capture
when it finally happens --- including pre-trigger history so you see what led up
to it. Add lightweight always-on instrumentation (RTT logging, a circular trace
buffer in RAM, error counters, min/max timestamps) so the fault leaves evidence.
For corruption, an MPU region or a watchpoint can trap the moment it occurs.
\trap the trap is trying to ``catch it live'' by staring or single-stepping ---
intermittent faults won't cooperate. The senior approach builds a \emph{trap}
that captures the event automatically (trigger + pre-trigger buffer + persistent
logging), so one occurrence yields the data. This is exactly how flaky
field/integration defects (the resume's HSI work) get cornered.

\question{Q3.4}{Design a bring-up plan for a new flight-control board so a fault
is localised, not mysterious.}
Bring it up in dependency order with a \emph{measurement and pass/fail at each
stage}, never skipping ahead: (1) power --- every rail's voltage, ripple, and
sequencing on the scope, current draw within budget; (2) clocks --- crystal/PLL
frequency confirmed on a clock-out pin; (3) reset/brownout --- clean release; (4)
debug access --- probe connects, IDCODE matches; (5) a minimal ``hello'' (GPIO
blink) proving the toolchain/startup/linker are right (Chapter 2); (6) each
peripheral and bus individually, scoped/decoded, with a known-good loopback or
reference device; (7) only then integration. Record every measurement against the
expected value so any deviation is caught at its stage. \trap the trap is flashing
the full application first and debugging top-down into a fog; the disciplined
bottom-up sequence means a failure at, say, the UART stage can't be confused with
a power or clock fault, because those already passed. That staged, instrumented,
documented bring-up \emph{is} the integration-and-V\&V rigour the resume points
to.

%=====================================================================
\section{Coding Gym}
%=====================================================================

\begin{partnote}
Four host-compilable problems (verified under
\code{gcc -Wall -Wextra -std=c11}) for the computational side of bench
debugging: decoding a captured waveform, SWD packet parity, a bring-up state
machine, and measuring frequency from samples --- the kind of small tools you
write to turn captures into answers.
\end{partnote}

%---------------------------------------------------------------
\begin{gymproblem}{Decode a UART byte from an oversampled capture}

\textbf{Problem.} Given a logic-analyser capture oversampled at a known rate,
recover the transmitted byte by sampling each bit centre.

\textbf{Hints.}
\begin{itemize}
\item The capture starts at the start-bit edge; data is LSB-first; sample at
  \code{os + os/2 + i*os}.
\end{itemize}

\attempt

\textbf{Solution.}
\begin{cgym}
#include <stdint.h>
#include <stdio.h>

static uint8_t uart_decode(const uint8_t *s, int os) {
    uint8_t b = 0;
    for (int i = 0; i < 8; i++) {
        int idx = os + os / 2 + i * os;          /* centre of data bit i (LSB first) */
        b = (uint8_t)(b | ((s[idx] & 1u) << i));
    }
    return b;
}

int main(void) {
    /* 'A' = 0x41, oversample 4: start, data LSB-first 1000 0010, stop, idle */
    int bits[11] = {0, 1,0,0,0,0,0,1,0, 1, 1};
    uint8_t s[44];
    for (int i = 0; i < 11; i++)
        for (int k = 0; k < 4; k++) s[i*4 + k] = (uint8_t)bits[i];
    printf("decoded=0x%02X\n", uart_decode(s, 4));   /* 0x41 */
    return 0;
}
\end{cgym}

This is what a logic analyser's UART decoder does internally and what you'd write
to post-process a raw capture: skip the start bit, sample each data bit at its
centre (most settled point), and assemble LSB-first. If your decode comes out
byte-reversed you sampled MSB-first; if every byte is wrong, the assumed baud
(bit spacing) doesn't match the capture --- the same diagnoses as Chapter 5, now
from the captured waveform.

\followups
\begin{enumerate}
\item \emph{``Add a parity check.''} $\rightarrow$ Sample the parity-bit centre
  and compare to the computed parity (Chapter 5).
\item \emph{``Detect a framing error in the capture.''} $\rightarrow$ Verify the
  stop-bit centre is high; if low, framing error.
\item \emph{``Auto-detect the baud from the capture.''} $\rightarrow$ Measure the
  narrowest pulse width $\rightarrow$ one bit time $\rightarrow$ baud.
\end{enumerate}
\end{gymproblem}

%---------------------------------------------------------------
\begin{gymproblem}{SWD packet parity}

\textbf{Problem.} Compute the even parity bit used in an SWD transfer's data
phase.

\textbf{Hints.}
\begin{itemize}
\item SWD uses even parity; XOR-reduce the data word's bits.
\end{itemize}

\attempt

\textbf{Solution.}
\begin{cgym}
#include <stdint.h>
#include <stdio.h>

static int swd_parity(uint32_t data) {
    int p = 0;
    while (data) { p ^= 1; data &= data - 1; }   /* count 1s mod 2 */
    return p;
}

int main(void) {
    printf("parity(0x5)=%d parity(0x7)=%d\n", swd_parity(0x5), swd_parity(0x7));
    /* 0x5 has two 1s -> 0 ; 0x7 has three -> 1 */
    return 0;
}
\end{cgym}

Each SWD packet protects its request and data phases with a parity bit; the host
and target each compute it and a mismatch flags a corrupted transfer (often a
signal-integrity problem on SWDIO/SWCLK --- too-long wires, missing series
resistor, or too-high clock). Understanding that the debug link itself is a
\emph{protocol with integrity checking} explains why an unreliable probe
connection throws ``parity'' or ``wire ACK'' errors --- and that the fix is
electrical (shorten/clean the wires, lower the SWD clock), the same instrument
discipline as any other bus.

\followups
\begin{enumerate}
\item \emph{``SWD ACK values?''} $\rightarrow$ OK (001), WAIT (010), FAULT (100)
  --- WAIT means retry, FAULT means a sticky error.
\item \emph{``Why do flaky SWD connections fail at high clock?''} $\rightarrow$
  Signal integrity on SWDIO degrades; lower the clock or shorten wires.
\item \emph{``JTAG integrity vs SWD?''} $\rightarrow$ JTAG has more wires but the
  same need for clean signals and reasonable TCK rate.
\end{enumerate}
\end{gymproblem}

%---------------------------------------------------------------
\begin{gymproblem}{Board bring-up state machine}

\textbf{Problem.} Model the bring-up sequence as a state machine that advances on
success and fails fast.

\textbf{Hints.}
\begin{itemize}
\item Power $\rightarrow$ clock $\rightarrow$ reset $\rightarrow$ JTAG
  $\rightarrow$ peripherals $\rightarrow$ done; any failure $\rightarrow$ fail.
\end{itemize}

\attempt

\textbf{Solution.}
\begin{cgym}
#include <stdio.h>

typedef enum { BU_POWER, BU_CLOCK, BU_RESET, BU_JTAG, BU_PERIPH, BU_DONE, BU_FAIL } bu_state_t;

static bu_state_t bu_next(bu_state_t s, int ok) {
    if (!ok) return BU_FAIL;                 /* fail fast: don't proceed on a bad step */
    switch (s) {
        case BU_POWER:  return BU_CLOCK;
        case BU_CLOCK:  return BU_RESET;
        case BU_RESET:  return BU_JTAG;
        case BU_JTAG:   return BU_PERIPH;
        case BU_PERIPH: return BU_DONE;
        default:        return s;
    }
}

int main(void) {
    bu_state_t st = BU_POWER; int steps = 0;
    while (st != BU_DONE && st != BU_FAIL) { st = bu_next(st, 1); steps++; }
    printf("end=%d steps=%d\n", st, steps);   /* end=5 (BU_DONE) steps=5 */
    return 0;
}
\end{cgym}

The point isn't the code --- it's the \emph{discipline} it encodes: each stage is
a precondition for the next, and a failure stops you \emph{at that stage} rather
than letting you chase a phantom firmware bug caused by a bad rail or dead clock
(Q3.4). Encoding bring-up as an explicit gated sequence (even on paper) is what
turns a mysterious dead board into a localised, one-line failure report. It's the
same fail-fast structure as a test sequence in Chapter 12.

\followups
\begin{enumerate}
\item \emph{``Attach a measurement to each state.''} $\rightarrow$ Rail voltage,
  clock frequency, IDCODE --- recorded pass/fail per stage.
\item \emph{``Why fail fast instead of continuing?''} $\rightarrow$ Later stages'
  results are meaningless if a precondition failed.
\item \emph{``Map to a real ATP.''} $\rightarrow$ Each state becomes a numbered
  test step with expected/actual (Chapter 12).
\end{enumerate}
\end{gymproblem}

%---------------------------------------------------------------
\begin{gymproblem}{Measure frequency from samples}

\textbf{Problem.} Compute a square wave's frequency from a sampled capture by
counting edges.

\textbf{Hints.}
\begin{itemize}
\item Count rising edges; frequency = edges / capture duration.
\end{itemize}

\attempt

\textbf{Solution.}
\begin{cgym}
#include <stdio.h>

static double measure_freq(const int *samp, int n, double sample_rate) {
    int edges = 0;
    for (int i = 1; i < n; i++)
        if (samp[i-1] == 0 && samp[i] == 1) edges++;   /* rising edges */
    double duration = (double)(n - 1) / sample_rate;
    return edges / duration;
}

int main(void) {
    int sq[16] = {0,1,0,1,0,1,0,1,0,1,0,1,0,1,0,1};
    printf("freq=%.0f Hz\n", measure_freq(sq, 16, 1000.0));   /* ~533 Hz */
    return 0;
}
\end{cgym}

Counting edges over a known time is how a scope/analyser derives frequency, and
how you'd verify a clock, PWM, or baud from a capture. The accuracy depends on the
sample rate (Nyquist: sample well above the signal frequency) and the capture
length (more cycles $\rightarrow$ better average). On the bench you'd confirm a
crystal/PLL or a timer's output this way during bring-up (Q2.3), turning ``is the
clock right?'' into a measured number rather than an assumption.

\followups
\begin{enumerate}
\item \emph{``Measure duty cycle too.''} $\rightarrow$ Fraction of samples high
  over a period --- verifies a PWM (Chapter 4).
\item \emph{``Why must the sample rate exceed 2$\times$ the signal?''}
  $\rightarrow$ Nyquist --- below it you alias and read a false frequency.
\item \emph{``Period from edges instead?''} $\rightarrow$ Average the
  sample-distance between successive rising edges; more precise for few cycles.
\end{enumerate}
\end{gymproblem}

%=====================================================================
\section{Link to Her Resume}
%=====================================================================

\begin{resumelink}
\textbf{Instrument-led root cause:} ``My root-cause method is instrument-led: I
scope the signal \emph{at the receiver} to separate electrical faults
(ringing $\rightarrow$ termination, slow edge $\rightarrow$ loading) from
logical ones, use a logic analyser to read the actual bytes on UART/SPI/CAN, and
a DMM for rails and continuity --- the oscilloscopes, DMMs, and analysers on my
resume.''

\medskip
\textbf{Debug probes \& flashing:} ``I flash and debug across boards with
ST-Link and J-Link over SWD/JTAG --- halting the core, setting breakpoints and
\emph{watchpoints} to catch memory corruption, and using RTT for low-overhead
logging from timing-critical code.''

\medskip
\textbf{Board bring-up:} ``I bring boards up in dependency order --- power,
clocks, reset, debug access, then peripherals one at a time --- with a
measurement and pass/fail at each stage, so a fault is localised rather than a
mystery. That's the same staged rigour as an ATP.''

\medskip
\small Maps directly to the resume: ST-Link and J-Link flashing and debug
across the autopilot and supporting boards, oscilloscope / DMM / logic-analyser
root-cause work, and multi-board bring-up at TASL. Deep project scripting is
Chapter 15, from \code{inputs/INTAKE.md} only.
\end{resumelink}

%=====================================================================
\section{Self-Test Scorecard}
%=====================================================================

\begin{scorecard}
\begin{enumerate}
\item Scope vs logic analyser --- which for an electrical fault, which for
  wrong data?
\item Why is triggering the key scope skill, and name two trigger types.
\item JTAG vs SWD --- pin counts and what both do.
\item Breakpoint vs watchpoint --- which catches ``who corrupted this
  variable?''
\item State the board bring-up order and the first step's tool.
\item Why probe at the receiver rather than the driver?
\item How do you measure interrupt latency on the bench?
\item What is RTT and why use it over a UART console?
\item Why does firmware sometimes work under the debugger but fail from reset?
\item How do you catch an intermittent fault you can't reproduce on demand?
\end{enumerate}
\end{scorecard}

\begin{answerkey}
\begin{enumerate}
\item Scope for electrical (edges/levels/ringing); logic analyser for wrong
  data/protocol content.
\item It captures the right moment stably; e.g.\ edge and serial/protocol (also
  pulse-width, single-shot).
\item JTAG 4--5 wires (TCK/TMS/TDI/TDO/nTRST), SWD 2 wires (SWCLK/SWDIO); both
  halt the core, access memory/registers, and flash.
\item A watchpoint (data breakpoint) --- it halts on access to the address,
  revealing the culprit instruction.
\item Power $\rightarrow$ clocks $\rightarrow$ reset $\rightarrow$ debug
  $\rightarrow$ peripherals; first step uses a DMM/scope on the rails.
\item The receiver sees the signal after cable/connector/termination effects ---
  where corruption actually occurs.
\item Toggle a GPIO at the event and another at ISR entry; trigger on the first,
  measure to the second (or use DWT cycles).
\item Real-Time Transfer --- streams logging through the debug probe via a RAM
  buffer; faster, no UART pin, minimal timing perturbation.
\item The debugger sets SP/clock/VTOR, halts the watchdog, initialises RAM, and
  serialises execution --- hiding startup/race/watchdog bugs that bite from cold
  reset.
\item Build a trap: trigger + pre-trigger capture on the fault signature, plus
  always-on logging (RTT/trace buffer/counters) or a watchpoint/MPU, so one
  occurrence yields data.
\end{enumerate}
\end{answerkey}

\begin{partnote}
\emph{End of Chapter 11.} Next: \textbf{V\&V, DO-178B \& Test Methodology} ---
white/black/grey box, coverage (statement/branch/MC/DC), requirements
traceability, ATP/QTP, HIL, fault injection, and the defect lifecycle --- the
verification framework around everything built so far.
\end{partnote}
